\section{The forced integer hierarchy}\label{sec:hierarchy}

\subsection{A counted pair and its exterior traces}

Let \(Y\) be \(F\) or a moving head \(G_N\) with \(N\ge2\), and set
\(R=(\partial_t+\partial_z^2+\beta_1)Y\).
On a fixed positively oriented circle \(\Gamma:|z-a|=r\) with
\(a\in\mathbb R\), assume a cleared collar and a count of two with
multiplicity. On compact portions of the chart take a buffer \(b>0\)
so the two occurrences lie in \(|z-a|\le r-b\). We use
\(\langle g\rangle=(2\pi i)^{-1}\oint_\Gamma g(z)\,dz\) and write
\begin{equation}\label{h:coordinates}
\begin{gathered}
 B=Y_z/Y,\qquad c=\tfrac12\langle zB\rangle,\qquad
 D=\tfrac12\langle z^2B\rangle-c^2,\qquad
 w=z-c,\qquad Q=w^2-D,\\
 A=Y/Q,\qquad U=A_z/A,\qquad f=R/Y .
\end{gathered}
\end{equation}
The argument principle makes \(c\) the mean and \(D\) the squared
half-difference of the two zeros. Both are real at real time.
The unit \(A\) extends holomorphically and is zero-free on the disk
and a collar. Each enclosed factor has modulus at least \(b\) on
\(\Gamma\), so \(|Q|\ge b^2\) there, also at \(D=0\). The
coordinates and the traces below are locally analytic in time for
the canonical fields, by the contour formulas and analyticity of
the source and finite moment construction.

For any rational weight \(k\) with poles only at the counted pair and
\(k(z)=O(|z|^{-2})\), put
\begin{equation}\label{h:trace}
 \tau[k]= -\langle kB\rangle=-\langle kU\rangle
        =\sum_{\rho\ \mathrm{exterior}} k(\rho).
\end{equation}
Every exterior zero is included with multiplicity. For the complete
field we use the even product and zero count \(O(x\log(2+x))\)
from \cite[Section~2]{agc}. Conjugating \(J_N\) by the alternating
diagonal sign matrix proves evenness of the head. The positive
spectral weights give \(|G_N(z)|\le e^{K_N|\operatorname{Im}z|}\), where
\(K_N=\sup_t\|J_N(t)\|\) on the compact time interval of the chart.
Jensen's formula at radius \(2x\) and \(G_N(0)=1\) give
\(n_N(x)\le 2K_Nx/\log 2\).
Hadamard factorization at order at most one pairs opposite factors
as \(1-z^2/\rho^2\). Evenness removes the possible linear exponential
and normalization fixes the constant. These products thus have locally
uniformly convergent logarithmic derivatives on the cleared collar
when summed in opposite pairs. Both zero counts give absolute
inverse-square summability. To prove \eqref{h:trace}, one exterior
pole contributes \(-k(\rho)\) to \(\langle kB\rangle\), by summing
the residues of \(k(z)/(z-\rho)\) over the counted poles and using
its vanishing residue at infinity. The rational target contribution
\(kQ_z/Q\) has all its poles inside \(\Gamma\) and zero total residue.
Any constant regularizer likewise integrates to zero.

The regular traces are
\begin{equation}\label{h:st}
 S_n=\tau[Q^{-n}]\quad(n\ge1),\qquad
 T_n=\tau[wQ^{-n}]\quad(n\ge2),\qquad M=S_1.
\end{equation}
Where \(D\ne0\), take a local square root \(s^2=D\) and put
\(z_\pm=c\pm s\), \(h=2s\), \(X=z-z_+\), and \(Z=z-z_-\).
The split traces, including both axes, are
\begin{equation}\label{h:split}
 k_{p,q}=X^{-p}Z^{-q},\qquad
 P_{p,q}=\tau[k_{p,q}],\qquad
 p,q\in\mathbb Z_{\ge0},\quad p+q\ge2 .
\end{equation}
All these weights have the required decay, as do their spatial and
coordinate derivatives on the respective charts.

\subsection{Motion of the pair and differentiation of a trace}

A dot denotes forward physical time. The normalized equation gives
\(Y_t/Y=-Y_{zz}/Y-\beta_1+f\) on the collar and hence
\begin{equation}\label{h:burgers}
 \dot B=-B_{zz}-2BB_z+f_z .
\end{equation}
Differentiating the counted moments gives
\(\dot c=-\langle Y_t/Y\rangle/2\) and
\(\dot D=-\langle wY_t/Y\rangle\). In \(Y_{zz}/Y\), write the
quadratic contribution as \(2/Q+2(Q_z/Q)U\) and the holomorphic
contribution as \(A_{zz}/A\). Since
\(\langle1/Q\rangle=0\), \(\langle w/Q\rangle=1\), and
\(wQ_z/Q=2+2D/Q\), the pair equations are
\begin{equation}\label{h:pair}
 \begin{gathered}
 \dot c=\langle (Q_z/Q)U\rangle+\gamma,\qquad
 \dot D=2-4DM+\delta,\\
 \gamma=-\tfrac12\langle f\rangle,\qquad
 \delta=-\langle wf\rangle .
 \end{gathered}
\end{equation}
The center term before \(\gamma\) is \(U(z_+)+U(z_-)\), with value
\(2U(c)\) at a double. At a double \(f=(R/A)/w^2\), so
\(\gamma=-\tfrac12(R/A)_z(c)\) and \(\delta=-(R/A)(c)\).
In particular \(\dot D=2\) at a double of the complete field.

\begin{proposition}[Full trace derivative]\label{h:forceprop}
For a kernel \(k(z;c,D)\) whose coordinate derivatives satisfy
the conditions in \eqref{h:trace}, the full derivative is
\begin{equation}\label{h:force}
 \begin{split}
 \frac{d}{dt}\tau[k]
   &=\langle k(B_{zz}+2BB_z)\rangle
               +\dot c\,\tau[k_c]+\dot D\,\tau[k_D]+\langle k_z f\rangle,\\
 \varepsilon[k]
   &=\langle k_z f\rangle+\gamma\tau[k_c]+\delta\tau[k_D]\\
   &=\big\langle\big(k_z-\tfrac12\tau[k_c]-w\tau[k_D]\big)f\big\rangle .
 \end{split}
\end{equation}
Here \(\varepsilon[k]\) is the part contributed by the residual
\(R\). For a holomorphic field variation with first derivative \(V\)
on the counted chart, the differentials of \(c,D\) and the contour
trace are \(\gamma,\delta,\varepsilon[k]\) with \(f\) replaced by \(V/Y\).
\end{proposition}
\begin{proof}
Differentiate \(-\langle kB\rangle\) on the fixed circle and insert
\eqref{h:burgers}. The identity
\(-\langle k f_z\rangle=\langle k_z f\rangle\) follows by
integration of \((kf)_z\). The first variation uses
\(dB=(V/Y)_z\) and the same integration by parts for the counted
moments. The formulas keep both changes of the kernel, even when
\(V/Y\) has a holomorphic part on the disk.
\end{proof}

For a head, \eqref{eq:finite-defect} shows that these residual terms are
\(\beta_N\beta_{N-1}\) times the derivatives in the last-link
direction, with the other links fixed. For \(F\), all three vanish.
The gap gradient \(-\langle wV/Y\rangle\) supplies the signed first
variation in Section~\ref{sec:signed}. The finite change in each
kernel enters the degree comparison in Section~\ref{sec:degree}.

\subsection{Split rows and axes}

When \(D\ne0\) the parts of \(\dot z_\pm\) contributed by \(R\) are
\(v_\pm=\gamma\pm\delta/h\), and
\begin{equation}\label{h:split-gap}
 \dot h=4/h-2hM+2\delta/h .
\end{equation}
For a split weight the residual part of \eqref{h:force} is
\begin{equation}\label{h:split-force}
 \begin{split}
 \varepsilon_{p,q}={}&p\big(v_+P_{p+1,q}-\langle k_{p+1,q}f\rangle\big)\\
                    &+q\big(v_-P_{p,q+1}-\langle k_{p,q+1}f\rangle\big).
 \end{split}
\end{equation}
A summand with vanishing prefactor is omitted. For \(m,n\ge1\) set
\[
 C_{m,n}=\sum_{i=2}^m P_{i,0}P_{m+1-i,n}
         +\sum_{j=1}^{n-1}P_{m,j}P_{0,n+1-j}.
\]
On the axes put
\(C_{m,0}=\sum_{i=2}^{m-1}P_{i,0}P_{m+1-i,0}\) and
\(C_{0,n}=\sum_{j=2}^{n-1}P_{0,j}P_{0,n+1-j}\).
An empty sum is zero.

\begin{theorem}[Forced split recurrence]\label{h:splitthm}
Write \(A_h=2/h-hM\). The mixed rows \((p,q\ge1)\) satisfy
\begin{equation}\label{h:mixed}
\begin{split}
 \dot P_{p,q}={}&p\{C_{p+1,q}-(p+3)P_{p+2,q}
                   -(q+2)P_{p+1,q+1}+A_hP_{p+1,q}\}\\
 &+q\{C_{p,q+1}-(p+2)P_{p+1,q+1}
                   -(q+3)P_{p,q+2}-A_hP_{p,q+1}\}
   +\varepsilon_{p,q}.
\end{split}
\end{equation}
The two axes have \(p\ge2\) and \(q\ge2\), respectively, and satisfy
\begin{equation}\label{h:axes}
 \begin{split}
 \dot P_{p,0}&=p\{C_{p+1,0}-(p+3)P_{p+2,0}-2P_{p+1,1}
                          +(2/h)P_{p+1,0}\}+\varepsilon_{p,0},\\
 \dot P_{0,q}&=q\{C_{0,q+1}-2P_{1,q+1}-(q+3)P_{0,q+2}
                          -(2/h)P_{0,q+1}\}+\varepsilon_{0,q}.
 \end{split}
\end{equation}
\end{theorem}

\begin{proof}
We give a contour derivation of the heat terms of
\cite[Eqs.~(3.5)--(3.7)]{agc} and the moving-head corrections of
\cite{jh02}. Set \(B_Q=Q_z/Q\) and
\(\mathcal I[k]=\langle kU^2\rangle\).
With \(B=B_Q+U\), move the derivatives in the heat term of
\eqref{h:force} onto \(k\) by integration by parts to obtain
\begin{equation}\label{h:heat-term}
 \langle k(B_{zz}+2BB_z)\rangle
     =-\tau[k_{zz}-2B_Qk_z]-\mathcal I[k_z].
\end{equation}
The terms containing only \(B_Q\) are rational, have all their
poles enclosed and decay at infinity with zero residue.

For the quadratic term, integrate the divided difference
\(-[k(u)-k(v)]/(u-v)\) against \(U(u)U(v)\), with \(u\) on
\(\Gamma\) and \(v\) on a smaller circle still enclosing the pair.
Both integrals are normalized by \((2\pi i)^{-1}\).
The \(k(u)\) term vanishes by the inner Cauchy integral. Integrating
the \(k(v)\) term first in \(u\) gives \(\mathcal I[k]\).
For \(k=k_{m,n}\), subtracting each reciprocal power gives
\begin{equation}\label{h:split-dd}
 -\frac{k_{m,n}(u)-k_{m,n}(v)}{u-v}
   =\sum_{i=1}^{m} k_{i,0}(u)k_{m+1-i,n}(v)
    +\sum_{j=1}^{n} k_{m,j}(u)k_{0,n+1-j}(v).
\end{equation}
Products integrate as products of the two linear functionals
\(-\langle kU\rangle\). For a first-power factor these functionals
are \(-U(z_+)\) or \(-U(z_-)\). Thus for mixed \(m,n\),
\(\mathcal I[k_{m,n}]=C_{m,n}-(U(z_+)+U(z_-))P_{m,n}\).
When \(m\ge3,n=0\) its two endpoint terms instead give
\(-2U(z_+)P_{m,0}\). The opposite axis gives
\(-2U(z_-)P_{0,n}\) for \(n\ge3\).

In \eqref{h:heat-term} use
\(\partial_z k_{p,q}=-p k_{p+1,q}-q k_{p,q+1}\).
In the block with weight \(l=k_{p+1,q}\), for example, the
expression in braces is
\(\tau[l_z-2B_Ql]+\mathcal I[l]+\dot z_+\tau[l]
 -\langle lf\rangle\).
The full velocity at a simple target is
\(\dot z_+=2/h+2U(z_+)+v_+\), and the other is
\(\dot z_-=-2/h+2U(z_-)+v_-\).
These follow from \eqref{h:pair} using
\(U(z_+)-U(z_-)=-hM\), itself the residue identity for
\(-\langle U/Q\rangle\). The two endpoint terms therefore give
\(\pm A_h\) in mixed blocks and \(\pm2/h\) on the axes.
The linear terms follow from \(B_Q=1/X+1/Z\).
This proves the formulas without differentiating any exterior
zero, and so also covers exterior multiplicities.
\end{proof}

\subsection{Regular rows through a collision}

\begin{theorem}[Forced symmetric recurrence]\label{h:regularthm}
In the ranges in \eqref{h:st}, the equations on the full counted
chart are
\begin{align}
 \dot S_n={}&2n\sum_{j=1}^n S_jS_{n+1-j}
   +2n\sum_{j=2}^n\big(T_jT_{n+2-j}+DS_jS_{n+2-j}\big)\notag\\
 &-4n(n+2)S_{n+1}-4n(n+3)DS_{n+2}+\varepsilon^S_n ,
 \label{h:S-row}\\
 \dot T_n={}&2(2n-1)\sum_{j=2}^n T_jS_{n+1-j}
    +4nD\sum_{j=2}^n T_jS_{n+2-j}\notag\\
 &-4(n^2+n-1)T_{n+1}-4n(n+3)DT_{n+2}+\varepsilon^T_n .
 \label{h:T-row}
\end{align}
The residual terms are
\begin{equation}\label{h:ST-force}
\begin{split}
 \varepsilon^S_n={}&-2n\langle wQ^{-n-1}f\rangle
                      +2n\gamma T_{n+1}+n\delta S_{n+1},\\
 \varepsilon^T_n={}&\langle ((1-2n)Q^{-n}-2nDQ^{-n-1})f\rangle\\
                  &+\gamma\big((2n-1)S_n+2nDS_{n+1}\big)
                      +n\delta T_{n+1}.
\end{split}
\end{equation}
\end{theorem}

\begin{proof}
Both kernels satisfy \(k_c=-k_z\). The even kernel has
\(k_z=-2nwQ^{-n-1}\), \(k_D=nQ^{-n-1}\), and the odd kernel has
\(k_z=(1-2n)Q^{-n}-2nDQ^{-n-1}\), \(k_D=nwQ^{-n-1}\).
This yields \eqref{h:ST-force}.

To use the nested-circle identity from the split proof, write
\(\alpha=Q(u)^{-1}\), \(\lambda=Q(v)^{-1}\) and
\(L_j=\sum_{i=0}^{j-1}\alpha^{j-1-i}\lambda^i\) with \(L_0=0\).
The negative divided differences of \(Q^{-j}\) and \(wQ^{-j}\)
are respectively
\[
 (w(u)+w(v))\alpha\lambda L_j,\qquad
 \alpha\lambda L_{j-1}+
   (D+w(u)w(v))\alpha\lambda L_j .
\]
These follow from \(Q(u)-Q(v)=(u-v)(w(u)+w(v))\).
For this computation only write \(T_1=-\langle wU/Q\rangle\).
Then
\begin{align*}
 \mathcal I[Q^{-j}]&=2\sum_{i=1}^{j} T_iS_{j+1-i},\\
 \mathcal I[wQ^{-j}]&=\sum_{i=1}^{j-1} S_iS_{j-i}
             +\sum_{i=1}^{j}(T_iT_{j+1-i}+DS_iS_{j+1-i}).
\end{align*}
The center speed before adding \(\gamma\) is \(-2T_1\). In the
even row this cancels the two \(T_1\) endpoints from
\(2n\mathcal I[wQ^{-n-1}]\). In the odd row it cancels
the \(T_1\) terms of
\((2n-1)\mathcal I[Q^{-n}]+2nD\mathcal I[Q^{-n-1}]\).
The \(-4DM\) part of the gap speed cancels the \(DS_1\) endpoint
products in the even row and the \(T_{n+1}S_1\) term in the
\(D\)-weighted odd sum. The remaining linear terms come from
\(-\tau[k_{zz}-2B_Q k_z]+2\tau[k_D]\). For \(Q^{-n}\) the first
operator inside \(\tau\), before the minus sign, has coefficients
\(2n(2n+5)\) on \(Q^{-n-1}\) and \(4n(n+3)D\) on \(Q^{-n-2}\).
For \(wQ^{-n}\) they are \(2(n+2)(2n-1)\) on \(wQ^{-n-1}\)
and \(4n(n+3)D\) on \(wQ^{-n-2}\).
Adding the \(2\tau[k_D]\) part gives exactly
\eqref{h:S-row}--\eqref{h:T-row}. The auxiliary \(T_1\) was a
contour value and cancels from the recurrences. Setting \(f=0\)
recovers \cite[Eqs.~(3.11)--(3.12)]{agc}.
\end{proof}

The lowest rows, which enter the degree and signed gap comparisons, are
\begin{equation}\label{h:low}
 \begin{split}
 \dot M &=2M^2-12S_2-16DS_3+\varepsilon^S_1,\\
 \dot S_2 &=8MS_2+4T_2^2+4DS_2^2-32S_3-40DS_4+\varepsilon^S_2,\\
 \dot T_2 &=(6M+8DS_2)T_2-20T_3-40DT_4+\varepsilon^T_2 .
 \end{split}
\end{equation}
Here \(\varepsilon^S_1=-2\langle wQ^{-2}f\rangle
+2\gamma T_2+\delta S_2\). In particular, comparing \(\dot M\)
requires the change of \(12S_2+16DS_3\) along with the change
of this entire residual term.

For \(D\ne0\), \(S_n=P_{n,n}\) and
\(T_n=(P_{n-1,n}+P_{n,n-1})/2\) for \(n\ge2\).
Conversely, for \(m=\max(p,q)\), the numerator of
\(k_{p,q}=(w-s)^{m-p}(w+s)^{m-q}/Q^m\) has degree at most
\(2m-2\) in \(w\). Reducing by \(w^2=Q+D\) expresses each split
trace in admissible \(S,T\) with coefficients polynomial in \(s\).
Near a double \(t_*\), this gives \(P_{p,q}=a_{p,q}(t)+s(t)b_{p,q}(t)\)
with analytic coefficients. Differentiating the kernel at \(s=0\)
gives
\[
 b_{p,q}(t_*)=(p-q)\sum_{\rho\ \mathrm{exterior}}
      (\rho(t_*)-c(t_*))^{-p-q-1}.
\]
If \(b_{p,q}(t_*)\dot D(t_*)\ne0\), the split derivative has a
term of order \(1/s\). The symmetric rows retain all higher
traces and stay regular at \(D=0\) by the contour buffer.
